\documentclass[11pt]{article}
\usepackage{setspace}
\usepackage{geometry}
\geometry{margin=1in}
\setstretch{1.1}

\begin{document}

\begin{flushright}
\textbf{Zhongli Wang} \\
Phone: +49 160 9541 6015 \\
Email: nkwangzhongli@gmail.com \\
Website: https://sites.google.com/view/zhongli-wang \\
Bielefeld, Germany \\
\end{flushright}

\vspace{0.25cm}

\noindent
\textbf{Selection Committee} \\
Environmental Change Institute \\
School of Geography and the Environment \\
University of Oxford \\
South Parks Road \\
Oxford, OX1 3QY \\

\vspace{0.25cm}

\noindent
\textbf{Application for Postdoctoral Researcher in the Political Economy of Net Zero}\\

\vspace{0.15cm}

\noindent Dear Dr Brown and Members of the Selection Committee,\\[0.1cm]

I am writing to apply for the position of Postdoctoral Researcher in the Political Economy of Net Zero at the Environmental Change Institute, University of Oxford. I am close to completing my Ph.D.\ in Economics at Bielefeld University and Universit\`a Cattolica del Sacro Cuore within the Marie Sk{\l}odowska--Curie EPOC programme. My research focuses on complex economic systems: how heterogeneous actors interact through networks, how firms respond to regulation and market constraints, and how institutional arrangements shape system-level outcomes. I am especially attracted to this role because it approaches net zero not only as a technical or financial target, but as a social, economic, and political transformation.

My doctoral dissertation, \emph{Strategic Behavior and Economic Dynamics in ICT Markets: An Agent-Based Framework}, develops computational and empirical approaches to studying market dynamics under regulation, consumer heterogeneity, network effects, congestion, and strategic firm investment. My empirical focus has been on digital and ICT markets, yet the analytical logic is closely connected to the political economy of net zero. In both settings, outcomes are shaped by interactions between firms, consumers, regulators, and institutional constraints. I have analyzed how micro-level decisions generate industry-level dynamics, how governance rules influence firm behavior, and how policy interventions may produce intended and unintended system-level effects.

A central strength I would bring to the project is my experience in agent-based modeling and complexity-informed economic analysis. In my job market paper, I study how firms make capacity-investment decisions when consumer demand is shaped by network effects and congestion. This required formalizing behavioral rules, running simulations, and interpreting the resulting market dynamics. These skills are relevant to net-zero sectors such as energy, housing, transport, and food, where investment decisions, consumer behavior, infrastructure constraints, and regulatory incentives interact in complex and sometimes path-dependent ways.

My work on data governance is also relevant to this project. In one paper, I examine how regulatory intensity affects firm competition, consumer welfare, and regional policy coordination. The broader question is how governance arrangements alter firm incentives and social outcomes. This connects closely with the project's focus on ownership structures, governance systems, business models, and value logics. For the political economy of net zero, I see these issues as central: who owns low-carbon assets, who controls investment decisions, who bears transition risks, and how different institutional arrangements affect public legitimacy.

I would conceptualize the UK's net-zero transition as a challenge of institutional design as much as technological substitution. The transition requires rapid decarbonization across various sectors including energy, housing, transport, and food. Policies in these sectors affect costs, mobility, employment, regional development, and public trust. For this reason, the political economy of net zero cannot be reduced to carbon prices or least-cost pathways. It must also examine ownership, accountability, affordability, coordination capacity, and the distribution of costs and benefits. A technically efficient transition may still fail politically if it is perceived as unfair, imposed, or captured by narrow private interests.

In this respect, I find Dr Brown's work on energy demand reduction especially relevant. Its emphasis on citizens' valuation of co-benefits, wellbeing, fairness, intergenerational concerns, and the limits of narrow cost-benefit assessment speaks directly to how I understand the challenge of net zero. Demand-side mitigation is not only a matter of reducing consumption; it also raises questions about quality of life, social justice, public acceptance, and the institutional conditions under which lifestyle changes are considered legitimate. I would be very interested in contributing to research that connects these concerns with a wider wellbeing-economy approach.

Beyond computational modeling, I bring experience in interdisciplinary and policy-oriented research. My coauthored work on \emph{Polytope Fraud Theory} connects financial regulation, accounting, machine learning, and investor protection, while another project on \emph{Agentic Workflows for Economic Research} explores how large-language-model agents can support transparent and reproducible research pipelines. These projects strengthened my ability to work across disciplinary boundaries, handle heterogeneous evidence, and translate technical analysis into broader institutional and policy implications. I would also be eager to further develop my qualitative research skills through interviews with policymakers, industry stakeholders, and other actors involved in net-zero transitions.

The Environmental Change Institute is an ideal environment for this work because of its commitment to interdisciplinary research on environmental change and its engagement with real-world policy challenges. I am particularly drawn to the opportunity to work with Dr Brown on a high-profile UKRI Future Leaders Fellowship project that aims to rethink the political economy of decarbonization. I would bring to the team a strong foundation in economic modeling, agent-based simulation, institutional analysis, and policy-oriented research, together with a collaborative attitude and a willingness to learn from other disciplinary traditions.

Thank you so much for considering my application! I would be delighted to discuss how my background in computational economics, complex systems, and institutional analysis can contribute to the project and to the wider research community at the Environmental Change Institute.

\vspace{0.35cm}

\noindent Sincerely,\\[0.95cm]
\noindent Zhongli Wang

\end{document}
